% GENERATED FILE. Edit canonical lessons, then run npm run generate:books.
% canonical-source-hash: fnv1a64:103e1fefe0699204
% canonical-lessons: ZH-C218-fengyi, ZH-W218-feng-wind, ZH-C218-nuanhuo, ZH-C218-menre, ZH-C218-ganzao, ZH-C218-chaoshi

\chapter{Wind and Warm Days}
\label{ch:zh-wind-and-warm-days}
\hlchaptermodality{218}

\begin{chapteropening}
\textbf{By the end of this chapter:} \emph{I can talk about a windbreaker, a mild day, muggy heat, dry air and damp.}

Complete the last lesson of chapter 218: I can talk about a windbreaker, a mild day, muggy heat, dry air and damp.
\end{chapteropening}

\section[\texorpdfstring{\zh{风衣} (fēngyī)}{风衣 (fēngyī)}]{\zh{风衣} (fēngyī) — a windbreaker; a light coat against the wind}

\label{lesson:ZH-C218-fengyi}



\begin{quote}
Before the new one: say the Mandarin for to ask around, then the Mandarin for to solve.
\end{quote}

\subsection*{You'll want to know: \zh{风衣}}
\textbf{\zh{风衣}} — \emph{fēngyī} — \textquotedblleft{}a windbreaker; a light coat against the wind\textquotedblright{}.

\begin{quote}
\emph{chūntiān fēng hěn dà wǒ chángcháng chuān fēngyī} — It's windy in spring, so I often wear a windbreaker.
\end{quote}

The new character, in pieces: \zh{风} — four strokes: a long stroke falling to the left, a bar that turns and hooks down to the right, and inside them a short stroke to the left and a dot. The next lesson puts it in your hand.

\begin{cousinweb}
\textbf{\zh{风衣}} (\emph{fēngyī}) is built from \zh{风} 'wind' and \zh{衣} 'garment', both inherited from Middle Chinese, \emph{pjuwng} and \emph{'j+j} (Guangyun, in Baxter's transcription). The dictionary explains the first, after an older book, as the messenger of heaven and earth.
\end{cousinweb}

\begin{grammarlens}[title={What you've built}]
One more word for describing things.
\end{grammarlens}

\subsection*{Your turn}
\begin{itemize}
\raggedright
  \item \emph{Say it:} \emph{fēngyī}
  \item \emph{Say it:} \emph{fēngyī}, once more
  \item \emph{Say it:} \emph{chūntiān fēng hěn dà wǒ chángcháng chuān fēngyī} — It's windy in spring, so I often wear a windbreaker.
  \item \emph{Recall:} say the Mandarin for to ask around, then the Mandarin for to solve, then say \emph{fēngyī} again
\end{itemize}

\begin{tcolorbox}[breakable,colback=teal!4,colframe=teal!35!black,title={Before you move on}]
What does \zh{风衣} mean? (\textquotedblleft{}A windbreaker; a light coat against the wind\textquotedblright{}.) Say it once more.
\end{tcolorbox}
\section[\texorpdfstring{\zh{风} (fēng)}{风 (fēng)}]{\zh{风} — a frame, and two short strokes inside}

\label{lesson:ZH-W218-feng-wind}



\begin{quote}
The word first, then the character inside it.

\begin{itemize}
\raggedright
  \item \emph{Recall:} say \emph{fēngyī} — \textbf{R1}, one lesson back, before it decays
\end{itemize}

\begin{quote}
\zh{衣} \zh{气}
\end{quote}

\emph{Read it:} each one aloud, with nothing said first — \emph{yī}, \emph{qì}
\end{quote}

\subsection*{Script — a frame, and two short strokes inside}
\hlblockfigure{\detokenize{figures/ZH-W218-feng-wind-filmstrip.pdf}}{How \zh{风} is written, stroke by stroke}

\begin{quote}
\zh{风}
\end{quote}

\textbf{Two pieces: a frame, then what sits inside it.} Count the pieces, not the strokes.

\textbf{The frame, two strokes.} A long stroke falling to the left, from the top down to the foot. Then a bar from its top that runs right, turns and falls to the lower right, flicking up at its end.

\textbf{Inside, two strokes.} A short stroke falling to the left, then a dot falling to the right across it, like a small cross.

2 + 2 = four.

\subsection*{Your turn}
\begin{itemize}
\raggedright
  \item \emph{Write it:} \zh{风}, all four, with the model in front of you
  \item \emph{Cover:} write it again from memory, and count the pieces, not the strokes
\end{itemize}

\begin{tcolorbox}[breakable,colback=teal!4,colframe=teal!35!black,title={Before you move on}]
How many strokes in \zh{风}? (\textbf{Four}.) What comes first? (\textbf{The long stroke falling to the left}.) What sits inside? (\textbf{A short fall to the left and a dot}.)

Source: \href{https://raw.githubusercontent.com/chanind/hanzi-writer-data/68d10a4b21150cae5e1ebbd223eed289cf32d90c/data/\%E9\%A3\%8E.json}{Hanzi Writer Data \zh{风}, pinned stroke paths and medians}.
\end{tcolorbox}
\section[nuǎnhuo]{nuǎnhuo — pleasantly warm; mild}

\label{lesson:ZH-C218-nuanhuo}



\begin{quote}
Before the new one: say the Mandarin for to solve, then the Mandarin for a windbreaker.
\end{quote}

\subsection*{You'll want to know: nuǎnhuo}
\emph{nuǎnhuo} — \textquotedblleft{}pleasantly warm; mild\textquotedblright{}.

\begin{quote}
\emph{jīntiān hěn nuǎnhuo} — It's lovely and warm today.
\end{quote}

Its characters come later; for now it is a word for the ear.

\begin{cousinweb}
\emph{nuǎnhuo} is built from \emph{nuǎn} 'warm' and \emph{huo} 'mild', both inherited from Middle Chinese, \emph{nwanX} and \emph{hwa} (Guangyun, in Baxter's transcription). The dictionary lists the first beside another character glossed 'warm', and glosses the second 'in harmony' and 'neither hard nor soft'; in the word the second loses its tone and its vowel changes.
\end{cousinweb}

\begin{grammarlens}[title={What you've built}]
One more word for describing things.
\end{grammarlens}

\subsection*{Your turn}
\begin{itemize}
\raggedright
  \item \emph{Say it:} \emph{nuǎnhuo}
  \item \emph{Say it:} \emph{nuǎnhuo}, once more
  \item \emph{Say it:} \emph{jīntiān hěn nuǎnhuo} — It's lovely and warm today.
  \item \emph{Recall:} say the Mandarin for to solve, then the Mandarin for a windbreaker, then say \emph{nuǎnhuo} again
\end{itemize}

\begin{tcolorbox}[breakable,colback=teal!4,colframe=teal!35!black,title={Before you move on}]
What does nuǎnhuo mean? (\textquotedblleft{}Pleasantly warm; mild\textquotedblright{}.) Say it once more.
\end{tcolorbox}
\section[mēnrè]{mēnrè — hot and stuffy; muggy}

\label{lesson:ZH-C218-menre}



\begin{quote}
Before the new one: say the Mandarin for a windbreaker, then the Mandarin for pleasantly warm.
\end{quote}

\subsection*{You'll want to know: mēnrè}
\emph{mēnrè} — \textquotedblleft{}hot and stuffy; muggy\textquotedblright{}.

\begin{quote}
\emph{jīntiān tài mēnrè le} — It's so muggy today.
\end{quote}

Its characters come later; for now it is a word for the ear.

\begin{cousinweb}
\emph{mēnrè} is built from \emph{mēn} 'stifled' and \emph{rè} 'hot', both inherited from Middle Chinese, \emph{mwonH} and \emph{nyet} (Guangyun, in Baxter's transcription), which glosses the first 'vexed' and explains the second as burning, as if by fire. The first is listed with a departing tone, so the level tone of mēn is not the regular outcome.
\end{cousinweb}

\begin{grammarlens}[title={What you've built}]
One more word for describing things.
\end{grammarlens}

\subsection*{Your turn}
\begin{itemize}
\raggedright
  \item \emph{Say it:} \emph{mēnrè}
  \item \emph{Say it:} \emph{mēnrè}, once more
  \item \emph{Say it:} \emph{jīntiān tài mēnrè le} — It's so muggy today.
  \item \emph{Recall:} say the Mandarin for a windbreaker, then the Mandarin for pleasantly warm, then say \emph{mēnrè} again
\end{itemize}

\begin{tcolorbox}[breakable,colback=teal!4,colframe=teal!35!black,title={Before you move on}]
What does mēnrè mean? (\textquotedblleft{}Hot and stuffy; muggy\textquotedblright{}.) Say it once more.
\end{tcolorbox}
\section[gānzào]{gānzào — dry air; dry (of weather or skin)}

\label{lesson:ZH-C218-ganzao}



\begin{quote}
Before the new one: say the Mandarin for pleasantly warm, then the Mandarin for hot and stuffy.
\end{quote}

\subsection*{You'll want to know: gānzào}
\emph{gānzào} — \textquotedblleft{}dry air; dry (of weather or skin)\textquotedblright{}.

\begin{quote}
\emph{dōngtiān zhèr hěn gānzào} — Winters here are very dry.
\end{quote}

Its characters come later; for now it is a word for the ear.

\begin{cousinweb}
\emph{gānzào} is built from \emph{gān} 'dry' and \emph{zào} 'parched'. The first is inherited from Middle Chinese \emph{kan} (Guangyun, in Baxter's transcription); the dictionary glosses the second, \emph{sawX} there, with this very word, 'dry'.
\end{cousinweb}

\begin{grammarlens}[title={What you've built}]
One more word for describing things.
\end{grammarlens}

\subsection*{Your turn}
\begin{itemize}
\raggedright
  \item \emph{Say it:} \emph{gānzào}
  \item \emph{Say it:} \emph{gānzào}, once more
  \item \emph{Say it:} \emph{dōngtiān zhèr hěn gānzào} — Winters here are very dry.
  \item \emph{Recall:} say the Mandarin for pleasantly warm, then the Mandarin for hot and stuffy, then say \emph{gānzào} again
\end{itemize}

\begin{tcolorbox}[breakable,colback=teal!4,colframe=teal!35!black,title={Before you move on}]
What does gānzào mean? (\textquotedblleft{}Dry air; dry (of weather or skin)\textquotedblright{}.) Say it once more.
\end{tcolorbox}
\section[cháoshī]{cháoshī — damp; humid}

\label{lesson:ZH-C218-chaoshi}



\begin{quote}
Before the new one: say the Mandarin for hot and stuffy, then the Mandarin for dry air.
\end{quote}

\subsection*{You'll want to know: cháoshī}
\emph{cháoshī} — \textquotedblleft{}damp; humid\textquotedblright{}.

\begin{quote}
\emph{xiàtiān zhèr hěn cháoshī} — Summers here are very humid.
\end{quote}

Its characters come later; for now it is a word for the ear.

\begin{cousinweb}
\emph{cháoshī} is built from \emph{cháo} 'tide' and \emph{shī} 'wet', both inherited from Middle Chinese, \emph{drjew} and \emph{syip} (Guangyun, in Baxter's transcription), which glosses the first 'tide water' and the second 'soaked with water'.
\end{cousinweb}

\begin{grammarlens}[title={What you've built}]
One more word for describing things.
\end{grammarlens}

\subsection*{Your turn}
\begin{itemize}
\raggedright
  \item \emph{Say it:} \emph{cháoshī}
  \item \emph{Say it:} \emph{cháoshī}, once more
  \item \emph{Say it:} \emph{xiàtiān zhèr hěn cháoshī} — Summers here are very humid.
  \item \emph{Recall:} say the Mandarin for hot and stuffy, then the Mandarin for dry air, then say \emph{cháoshī} again
\end{itemize}

\begin{tcolorbox}[breakable,colback=teal!4,colframe=teal!35!black,title={Before you move on}]
What does cháoshī mean? (\textquotedblleft{}Damp; humid\textquotedblright{}.) Say it once more.
\end{tcolorbox}
